\documentclass[aspectratio=169,11pt]{beamer}

% ============================================
% PACKAGES
% ============================================
\usepackage[utf8]{inputenc}
\usepackage[T1]{fontenc}
\usepackage[ngerman]{babel}
\usepackage{lmodern}
\usepackage{tcolorbox}
\tcbuselibrary{skins, hooks}
\usepackage{listings}
\usepackage{xcolor}
\usepackage{fontawesome5}
\usepackage{tikz}
\usetikzlibrary{trees, shapes, positioning}

% Modern Font (if installed, otherwise falls back smoothly)
\IfFileExists{FiraSans.sty}{\usepackage[sfdefault]{FiraSans}}{}
\renewcommand{\familydefault}{\sfdefault}

% ============================================
% COLORS & DESIGN SYSTEM
% ============================================
\definecolor{BgColor}{HTML}{FAFAFA}
\definecolor{Primary}{HTML}{2B3A42}
\definecolor{Accent}{HTML}{E84A5F}
\definecolor{Secondary}{HTML}{3F5765}
\definecolor{CodeBg}{HTML}{F0F0F0}
\definecolor{Success}{HTML}{28A745}

\setbeamercolor{background canvas}{bg=BgColor}
\setbeamercolor{title}{fg=Primary}
\setbeamercolor{frametitle}{bg=Primary, fg=white}
\setbeamercolor{structure}{fg=Accent}
\setbeamercolor{normal text}{fg=Primary}
\setbeamercolor{itemize item}{fg=Accent}
\setbeamercolor{itemize subitem}{fg=Secondary}

\setbeamertemplate{navigation symbols}{}
\setbeamertemplate{footline}{%
  \leavevmode%
  \hbox{%
  \begin{beamercolorbox}[wd=.333333\paperwidth,ht=2.25ex,dp=1ex,center]{author in head/foot}%
    \usebeamerfont{author in head/foot}\tiny Falko, Malte, Max, Jana
  \end{beamercolorbox}%
  \begin{beamercolorbox}[wd=.333333\paperwidth,ht=2.25ex,dp=1ex,center]{title in head/foot}%
    \usebeamerfont{title in head/foot}\tiny Grundlagen der Praktischen Informatik
  \end{beamercolorbox}%
  \begin{beamercolorbox}[wd=.333333\paperwidth,ht=2.25ex,dp=1ex,right]{date in head/foot}%
    \usebeamerfont{date in head/foot}\insertframenumber{} / \inserttotalframenumber\hspace*{2ex} 
  \end{beamercolorbox}}%
  \vskip0pt%
}

% ============================================
% CUSTOM BOXES
% ============================================
\tcbset{
    modernbox/.style={
        enhanced,
        boxrule=0pt,
        frame hidden,
        colback=white,
        coltitle=white,
        colbacktitle=Secondary,
        fonttitle=\bfseries,
        drop lifted shadow,
        arc=4pt,
        toptitle=2mm,
        bottomtitle=2mm,
        left=2mm,
        right=2mm
    },
    alertbox/.style={
        modernbox,
        colbacktitle=Accent
    },
    successbox/.style={
        modernbox,
        colbacktitle=Success
    }
}

% ============================================
% CODE LISTINGS
% ============================================
\lstdefinelanguage{Haskell}{
  keywords={class, data, deriving, do, else, if, in, import, let, module, of, then, type, where, instance},
  keywordstyle=\color{Accent}\bfseries,
  ndkeywords={print, main, return, Just, Nothing, Left, Right},
  ndkeywordstyle=\color{Secondary}\bfseries,
  identifierstyle=\color{Primary},
  sensitive=true,
  comment=[l]{--},
  commentstyle=\color{gray}\ttfamily\itshape,
  stringstyle=\color{teal}\ttfamily,
  morestring=[b]"
}

\lstset{
    language=Haskell,
    backgroundcolor=\color{CodeBg},
    basicstyle=\ttfamily\scriptsize,
    breaklines=true,
    frame=single,
    rulecolor=\color{CodeBg},
    frameround=tttt,
    numbers=left,
    numberstyle=\tiny\color{gray},
    xleftmargin=15pt,
    tabsize=4,
    showstringspaces=false
}

% ============================================
% TITLE PAGE SETTINGS
% ============================================
\title{\textbf{Tutorium 3}}
\subtitle{Typklassen \& Algebraische Datentypen}
\institute{Grundlagen der Praktischen Informatik}
\author{}
\date{\today}

\begin{document}

\begin{frame}[plain]
    \titlepage
\end{frame}

\begin{frame}{Inhaltsverzeichnis}
    \tableofcontents
\end{frame}

% ============================================
% THEORIE: TYPKLASSEN
% ============================================
\section{Typklassen}

\begin{frame}[fragile]{Warum Typklassen?}
    Ohne Typklassen müssen wir für jeden primitiven Datentyp (z.\,B. \texttt{Integer}, \texttt{Double}) eine eigene Funktion schreiben, auch wenn die Logik identisch ist.
    
    \begin{tcolorbox}[alertbox, title=Das Problem]
\begin{lstlisting}
greaterFive :: Integer -> Bool
greaterFive x = x > 5

greaterFive' :: Double -> Bool
greaterFive' x = x > 5
\end{lstlisting}
    \end{tcolorbox}
    \textbf{Lösung:} Wir verallgemeinern die Funktion mit \textbf{Typklassen} (vgl. Interfaces in Java).
\end{frame}

\begin{frame}[fragile]{Verwendung von Typklassen (\texttt{Num}, \texttt{Eq}, \texttt{Ord})}
    \begin{tcolorbox}[successbox, title=Die Lösung mit Typklassen]
\begin{lstlisting}
-- Erlaubt alle numerischen Typen, die ordnerbar sind
greaterFive'' :: (Ord a, Num a) => a -> Bool
greaterFive'' x = x > 5

numberToString :: (Eq a, Num a) => a -> String
numberToString 0 = "null"
numberToString _ = "andere Zahl"
\end{lstlisting}
    \end{tcolorbox}
    
    \begin{itemize}
        \item \texttt{Num}: Numerische Operationen (z.\,B. \texttt{Integer, Double, Float}).
        \item \texttt{Eq}: Erlaubt Überprüfung auf Gleichheit (\texttt{==}).
        \item \texttt{Ord}: Erlaubt Vergleiche (\texttt{>}, \texttt{<}). \textit{Wichtig: Nicht alles, was \texttt{Eq} ist, ist auch \texttt{Ord} (z.\,B. Komplexe Zahlen).}
    \end{itemize}
\end{frame}

% ============================================
% THEORIE: ALGEBRAISCHE DATENTYPEN
% ============================================
\section{Algebraische Datentypen (ADT)}

\begin{frame}[fragile]{Aufzählungsdatentypen (Summen-Typen)}
    Analog zu \texttt{enum} in Java zählen wir eine Menge von exakten Werten auf.
    
    \begin{tcolorbox}[modernbox, title=Beispiel: Wochentag]
\begin{lstlisting}
data Wochentag = Montag | Dienstag | Mittwoch | Donnerstag 
               | Freitag | Samstag | Sonntag
  deriving (Show, Eq)

printWochentagZahl :: Wochentag -> Integer
printWochentagZahl Montag = 1
printWochentagZahl _      = 0
\end{lstlisting}
    \end{tcolorbox}
    \small \textbf{Info:} \texttt{deriving (Show, Eq)} sorgt dafür, dass Haskell die Typen automatisch in der Konsole ausgeben (\texttt{Show}) und vergleichen (\texttt{Eq}) kann.
\end{frame}

\begin{frame}[fragile]{Produktdatentypen}
    Kombinieren verschiedene Werte zu einem neuen komplexen Typ (vgl. Klassen in Java).
    
    \begin{tcolorbox}[modernbox, title=Beispiel: Person]
\begin{lstlisting}
data Person = Person String Integer 
  deriving (Show, Eq)

printName :: Person -> String
printName (Person name alter) = name
\end{lstlisting}
    \end{tcolorbox}
\end{frame}

\begin{frame}[fragile]{Rekursive Datentypen (Visualisiert)}
    Ein ADT kann sich selbst referenzieren, um komplexe Strukturen aufzubauen, z.\,B. einen Baum.
    
    \begin{columns}[T]
        \begin{column}{0.48\textwidth}
            \begin{tcolorbox}[modernbox, title=Beispiel: Baum (Tree)]
\begin{lstlisting}
data Tree = Leaf Integer 
          | Node Tree Tree
\end{lstlisting}
            \end{tcolorbox}
            \small
            \textbf{Instanz:} \\
            \texttt{Node (Leaf 1) \\ \quad(Node (Leaf 2) (Leaf 3))}
        \end{column}
        \begin{column}{0.48\textwidth}
            \textbf{Speicherdarstellung (Abstrakt)}
            \vspace{0.3cm}
            \begin{center}
            \begin{tikzpicture}[level distance=1.2cm, sibling distance=2cm, every node/.style={circle, draw=Primary, fill=BgColor, minimum size=0.6cm, font=\small}]
                \node[rectangle, rounded corners=2pt] {Node}
                    child {node[fill=CodeBg] {1}}
                    child {node[rectangle, rounded corners=2pt] {Node}
                        child {node[fill=CodeBg] {2}}
                        child {node[fill=CodeBg] {3}}
                    };
            \end{tikzpicture}
            \end{center}
        \end{column}
    \end{columns}
\end{frame}

\begin{frame}[fragile]{Vordefinierte Typen: \texttt{Maybe} und \texttt{Either}}
    \begin{columns}[T]
        \begin{column}{0.48\textwidth}
            \begin{tcolorbox}[modernbox, title=\texttt{Maybe a}]
\begin{lstlisting}
data Maybe a = Nothing | Just a
\end{lstlisting}
            \end{tcolorbox}
            \small Verhindert Exceptions, indem ein Wert \texttt{Just x} oder eben \texttt{Nothing} zurückgegeben wird.
        \end{column}
        \begin{column}{0.48\textwidth}
            \begin{tcolorbox}[modernbox, title=\texttt{Either a b}]
\begin{lstlisting}
data Either a b = Left a | Right b
\end{lstlisting}
            \end{tcolorbox}
            \small Kann für zwei verschiedene Rückgabetypen verwendet werden (oft: \texttt{Left} = Error, \texttt{Right} = Resultat).
        \end{column}
    \end{columns}
\end{frame}

% ============================================
% LOGIK
% ============================================
\section{Exkurs: Aussagenlogik}

\begin{frame}{Wertetabellen \& Äquivalenz}
    \begin{itemize}
        \item Logische Ausdrücke lassen sich in \textbf{Wertetabellen} auswerten.
        \item Bei 3 Eingabevariablen (\texttt{a, b, c}) hat die Tabelle exakt $2^3 = 8$ Zeilen.
        \item \textbf{Äquivalenz:} $A \Leftrightarrow B$ ist genau dann wahr, wenn beide Wahrheitswerte übereinstimmen:
    \end{itemize}
    \vspace{0.3cm}
    \begin{center}
        \texttt{a <=> b} \quad ist äquivalent zu \quad \texttt{(a \&\& b) || (not a \&\& not b)}
    \end{center}
\end{frame}

% ============================================
% LIVE DEMO
% ============================================
\section{Praxis}

\begin{frame}[plain]
    \begin{center}
        \Huge \textbf{Live Demo} \\
        \vspace{0.5cm}
        \large \faLaptopCode~ \texttt{typklassen.hs} \& \texttt{algebraischeDatentypen.hs} \\
        \vspace{0.3cm}
        \normalsize Wir programmieren Typklassen und eigene ADTs zusammen im Code!
    \end{center}
    
    \vspace{0.5cm}
    \textbf{Aufgaben:}
    \begin{enumerate}
        \item \texttt{numberToMaybeString} (Fehlervermeidung durch \texttt{Maybe})
        \item Gerade Zahlen in \texttt{Left}, ungerade in \texttt{Right} packen (\texttt{Either})
        \item Logikrätsel gemeinsam durchgehen
    \end{enumerate}
\end{frame}

\end{document}
